\documentclass[11pt,a4paper]{article}
\usepackage[margin=25mm]{geometry}
\usepackage{amsmath,amsthm,amssymb,mathtools,booktabs,microtype}
\usepackage[unicode,colorlinks=true,linkcolor=blue!50!black,citecolor=blue!50!black,urlcolor=blue!60!black]{hyperref}
\usepackage{xcolor}
\hypersetup{pdftitle={Sharp rates for fixed-weight Ritz sweeps on p+1 spectral nodes},pdfauthor={}}
\newtheorem{theorem}{Theorem}
\newtheorem{lemma}[theorem]{Lemma}
\DeclareMathOperator{\osc}{osc}
\newcommand{\R}{\mathbb R}
\newcommand{\abs}[1]{\lvert #1\rvert}
\newcommand{\norm}[1]{\lVert #1\rVert}
\title{Exact Convergence Rates of Fixed-Weight LMSD\\ on \(p+1\) Eigenvalues}
\author{}
\date{25 November 2026 \quad\textbullet\quad Version 0.2.0}
\begin{document}
\maketitle
\begin{center}
\small AI-assisted research working paper. Author review is pending.\\
The release is anonymous and draws on joint research.
\end{center}
\begin{abstract}
A sweep of limited-memory steepest descent takes several steps with Ritz values computed from an earlier gradient history. We restrict the matrix to \(p+1\) distinct active eigenvalues and use memory \(p\) with one fixed positive spectral weight. On this support, a Lagrange formula identifies the unique balanced spectral energy and the sharp uniform gradient factor. The same interpolation formula gives a linear recurrence for the logarithms of the component ratios. Outside the balanced pair, a two-sweep estimate converts the growth of these ratios into doubly exponential decay in the number of completed sweeps. The proof also gives a uniform prefactor at the sharp factor and an exact three-node example.
\end{abstract}

\section{The delayed sweep}

Let \(H\) be a real symmetric positive definite matrix and consider the gradient of \(\tfrac12 x^\top Hx-b^\top x\). A gradient step with size \(s\) maps \(g\) to \((I-sH)g\). We work in the invariant subspace spanned by eigenspaces for
\[
  J=\{\lambda_0<\cdots<\lambda_p\},\quad \lambda_0>0,\quad p\ge1.
\]
Write \(E_i\) for the associated spectral projections. The eigenspaces may have arbitrary multiplicity. Define \(W=\sum_i\omega_i E_i\), where \(\omega_i>0\) is fixed. A Ritz value here is an eigenvalue of the compression of \(H\) to
\[
  \mathcal K_p(H,u)=\operatorname{span}\{u,Hu,\ldots,H^{p-1}u\}
\]
in the inner product \(\langle x,y\rangle_W=x^\top Wy\). The same \(W\) is used in every sweep. Since \(W\) acts by a scalar on each eigenspace, \(W^{1/2}\) commutes with \(H\). Thus for any polynomial \(P\),
\[
 W^{1/2}P(H)u=P(H)W^{1/2}u.
\]
The weighted projection is equivalent to Euclidean projection after changing variables to \(W^{1/2}u\). Use the actual dimension \(d\le p\) and write
\[
  P_u(t)=\prod_{j=1}^{d}(1-t/\theta_j(u)).
\]
For \(u=0\) define \(P_0=1\).

Let \(u_k\) and \(u_{k+1}\) be the gradients at the start and end of a completed sweep. The stored history consists of the gradients at which the steps were taken, so it does not include \(u_{k+1}\). In a sweep of \(p\) steps these gradients are polynomials in \(H\) applied to \(u_k\). Their degrees are \(0,\ldots,p-1\), and their leading coefficients are nonzero because all step sizes are positive. The triangular relation with \(u_k,Hu_k,\ldots,H^{p-1}u_k\) gives the same span. The Ritz values from this history determine the \emph{next} sweep, which starts at \(u_{k+1}\). Consequently
\begin{equation}\label{eq:rule}
 u_{k+2}=P_{u_k}(H)u_{k+1}.
\end{equation}
The order of the factors can change the gradients within a sweep. The factors commute, so their order does not change its final gradient.

We call the initial pair compatible if \(u_1=q(H)u_0\), with \(q\) a product of at most \(p\) factors \(1-t/\theta\), \(\theta\in[\lambda_0,\lambda_p]\). This assumption specifies one way to obtain \(u_1\) from \(u_0\), for example by a fixed-step warm-up with reciprocal step size in the given interval.

A distinct eigenvalue \(\lambda_i\) is active for \(u\) if \(E_i u\ne0\). Full support means that all \(p+1\) values are active. Multiplication by a polynomial in \(H\) cannot create a component that is zero, so the active sets are nested. If \(u_k\) has at most \(p\) active eigenvalues, its Krylov space contains each nonzero \(E_i u_k\): interpolation gives a polynomial of degree less than \(p\) selecting that component. The Ritz values are then exactly the active eigenvalues. Since \(u_{k+1}\) has no additional active components, \(P_{u_k}(H)u_{k+1}=0\). If only \(u_{k+1}\) has lost support, the same argument at the next index yields \(u_{k+3}=0\). After termination all gradients are set to zero.

Fletcher introduced the limited-memory sweep \cite{Fletcher}, and Curtis and Guo proved \(R\)-linear convergence on positive definite quadratics under a full-rank, uniformly conditioned history assumption \cite{CurtisGuo}. Rank compression for weighted sweeps on an arbitrary finite spectrum is treated by Yang and Yuan \cite{YangYuan}; a \href{https://taitai2025.github.io/assets/documents/research/lmsd-sweep-proof-note.pdf}{companion proof note} develops that framework in detail. Gu studies rolling and sampled historical Ritz rules, including nodal formulas on minimal supports and superlinear rates under reuse \cite[Section~3]{Gu}. Here the setting is one full delayed sweep per refresh with a fixed weight. Our purpose is to present the balanced attaining orbit, the uniform bound at its exact factor, and a proof based only on finite-dimensional identities on this support.

\section{The balanced polynomial and the logarithmic dynamics}

Set
\[
 D_i=\prod_{j\ne i}(\lambda_i-\lambda_j),\quad
 L_i(t)=\prod_{j\ne i}\frac{t-\lambda_j}{\lambda_i-\lambda_j},\quad
 c_i=\abs{L_i(0)},\quad
 \tau=(\textstyle\sum_i c_i)^{-1},\quad
 w_i=\tau c_i .
\]
The signs of \(L_i(0)\) are \((-1)^i\). The numbers \(w_i\) depend only on the eigenvalues and set the mean used in the proof. The prescribed \(\omega_i\) specify the algorithm. The \(w_i\) sum to one, and \(0<\tau<1\) since the alternating sum \(\sum_i L_i(0)\) is one. For any polynomial \(P\) of degree at most \(p\) with \(P(0)=1\),
\[
  1=\abs{\textstyle\sum_iL_i(0)P(\lambda_i)}
    \le\tau^{-1}\max_i\abs{P(\lambda_i)} .
\]
The unique minimizer of the nodal maximum is \(P_\star=\tau\sum_i(-1)^iL_i\). Equality forces \(P(\lambda_i)=(-1)^i\tau\) at every node. This finite-set interpolation problem is a special case of extremal growth at an external evaluation point \cite{Bos}. Its role here is to give the filter of a trajectory that sustains equality under the delayed update.

For a full-support anchor \(u\), put \(m_i(u)=\omega_i\norm{E_i u}^2\). The polynomial \(P_u\) is characterized by
\begin{equation}\label{eq:ortho}
   \sum_i m_i(u) P_u(\lambda_i)\lambda_i^j=0,\qquad 0\le j<p,
   \qquad P_u(0)=1 .
\end{equation}
To see this, let \(A\) be the orthogonal compression of \(H\) in the \(W\)-inner product. For \(0\le j\le p\), \(A^ju\) is the projection of \(H^ju\) onto the Krylov space. This follows by induction since the earlier powers already belong to the space. Cayley--Hamilton for \(A\) yields \eqref{eq:ortho}. The rectangular Vandermonde matrix in \eqref{eq:ortho} has full row rank, and its kernel is spanned by \(1/D_i\), because the coefficient of \(t^p\) on the two sides of \(t^j=\sum_i\lambda_i^jL_i(t)\) gives \(\sum_i \lambda_i^j/D_i=0\) for \(j<p\). Thus \(m_i(u)P_u(\lambda_i)\) is a scalar multiple of \(1/D_i\). Its scale is determined by evaluation at zero. The nodal values are nonzero and alternate in sign, so the roots lie strictly one to each gap. In particular a full-support pair stays full support.

Define the centered coordinates and the weighted exponential sum by
\begin{align*}
  Y_i(u)&=\log(m_i(u)\abs{D_i})
         -\sum_j w_j\log(m_j(u)\abs{D_j}),\\
  M(y)&=\sum_i w_i e^{y_i}.
\end{align*}
The vector \(Y(u)=0\) if and only if \(m_i(u)\) is proportional to \(1/\abs{D_i}\). We call this the balanced energy shape. Equations \eqref{eq:ortho} and the normalization give
\begin{equation}\label{eq:nodal}
 P_u(\lambda_i)=(-1)^i \tau\frac{e^{-Y_i(u)}}{M(-Y(u))}.
\end{equation}
Taking the logarithm of squared component sizes in \eqref{eq:rule} and subtracting the \(w\)-mean cancels the scalar denominator in \eqref{eq:nodal}. With \(Y_k=Y(u_k)\),
\begin{equation}\label{eq:recurrence}
 Y_{k+2}=Y_{k+1}-2Y_k .
\end{equation}
The quadratic form
\[
 Q(y,z)=4\norm y_2^2-2\langle y,z\rangle+2\norm z_2^2
       =\tfrac72\norm y_2^2+2\norm{z-y/2}_2^2
\]
is positive definite. Expanding \(Q(z,z-2y)\) gives \(2Q(y,z)\). By \eqref{eq:recurrence},
\[
 Q(Y_k,Y_{k+1})=2^kQ(Y_0,Y_1).
\]

At balance the filter \eqref{eq:nodal} is \(P_\star\). It changes the signs of the components and multiplies their absolute sizes by \(\tau\). Taking a balanced \(u_0\) and \(u_1=P_\star(H)u_0\) gives a compatible orbit that attains the factor \(\tau\) on all sweeps. Every full-support pair with both energy shapes balanced has the same norm ratio from \(u_1\) onward.

\section{From changing spectral shapes to norm decay}

Measure the gradient in the norm \(\norm u_*=\norm{H^{-1/2}W^{1/2}u}_2\), where \(\norm u_*^2=\sum_i m_i(u)/\lambda_i\). On this fixed spectrum the norm is equivalent to the Euclidean norm. The proportion of its square at node \(i\) is
\[
  \pi_{k,i}=\frac{m_i(u_k)/\lambda_i}{\sum_j m_j(u_k)/\lambda_j}
           =\frac{w_i e^{Y_{k,i}}}{M(Y_k)} .
\]
Here \(c_i=(\prod_j\lambda_j)/(\lambda_i\abs{D_i})\). Write \(x_k=\norm{u_{k+2}}_*^2/\norm{u_{k+1}}_*^2\). Equations \eqref{eq:nodal} and \eqref{eq:recurrence} give
\begin{equation}\label{eq:factor}
 x_k=\tau^2\frac{M(Y_{k+2})}{M(Y_{k+1})M(-Y_k)^2}.
\end{equation}
Centering matters: Jensen's inequality yields \(M(-Y_k)\ge1\). This will control the initial stay near the balanced shape.

For the large-coordinate region, set \(a=\min_i w_i>0\). For arbitrary vectors \(v,z\),
\[
 M(v+z)\le a^{-1}M(v)M(z).
\]
The product on the right contains the diagonal terms, each with coefficient \(w_i^2\ge aw_i\). Since \(Y_{k+2}=-Y_k-2Y_{k-1}\) for \(k\ge1\), multiplying two instances of \eqref{eq:factor} gives
\begin{align}
 x_{k-1}x_k
 &\le \frac{\tau^4 a^{-2}}{M(Y_k)M(-Y_k)}
 \le \tau^4 a^{-4}e^{-\osc(Y_k)} . \label{eq:two}
\end{align}
For the last inequality, take the two terms corresponding to the largest and smallest coordinates in the product of sums.

Every step factor in \eqref{eq:rule} is bounded in absolute value by \(K=\max(1,\lambda_p/\lambda_0-1)\). In particular \(x_k\le K^{2p}\). Each centered vector has \(\norm y_2\le \sqrt{p+1}\,\osc(y)\). The largest eigenvalue of the matrix of \(Q\) on each coordinate pair is \(3+\sqrt2\). Hence for
\[
 d=\frac{1}{2\sqrt{(p+1)(3+\sqrt2)}}\quad\text{and}\quad E_0=Q(Y_0,Y_1),
\]
at least one of \(\osc(Y_{j+1})\) and \(\osc(Y_{j+2})\) is at least \(d\,2^{(j+1)/2}\sqrt{E_0}\). Applying \eqref{eq:two} to the corresponding pair of factors in a triple gives
\begin{equation}\label{eq:triple}
 x_jx_{j+1}x_{j+2}
 \le K^{2p}\tau^4a^{-4}\exp(-d\,2^{(j+1)/2}\sqrt{E_0}) .
\end{equation}

\begin{theorem}[Sharp factor and uniform prefactor]\label{thm:sharp}
There is a constant \(C(J,W)<\infty\) such that every compatible starting pair satisfies
\[
 \norm{u_{n+1}}_*\le C(J,W)\tau^n(\norm{u_0}_*+\norm{u_1}_*)\qquad(n\ge0).
\]
The worst asymptotic root factor is \(\tau\) per full sweep and \(\tau^{1/p}\) per individual step. On any nonbalanced full-support trajectory, for \(n\ge3\),
\[
  \norm{u_{n+1}}_*
  \le K^{p(n-2)}\tau^2a^{-2}
       \exp\!\left(-\tfrac d2\,2^{(n-2)/2}\sqrt{E_0}\right)\norm{u_1}_* .
\]
Such a trajectory has root factor zero. The theorem makes no assumption on the multiplicities of the \(p+1\) distinct nodes.
\end{theorem}

\begin{proof}
For the last bound, apply \eqref{eq:triple} with \(j=n-3\), bound the preceding \(n-3\) squared factors by \(K^{2p}\), and take square roots. This proves double exponential decay up to an exponential prefactor whenever \(E_0>0\).

For the uniform bound, choose \(A>0\) so large that the right side of \eqref{eq:triple} is at most \(\tau^6\) whenever \(2^{j+1}E_0\ge A\). If \(E_0\ge A\), apply the bound to consecutive groups of three factors. At most two factors remain, each at most \(K^{2p}\). Relative to \(\tau^{2n}\), their contribution is bounded by \(\max\{1,(K^{2p}/\tau^2)^2\}\).

If \(0<E_0<A\), let \(N\) be the first index with \(2^NE_0\ge A\). For every \(m\le N\) we have \(Q(Y_m,Y_{m+1})<2A\). Positive definiteness of \(Q\) bounds \(M(Y_{m+1})\) above by a constant \(B_A\) independent of the initial pair and \(N\). Multiplication of \eqref{eq:factor} gives
\[
 \prod_{k=0}^{m-1}x_k
 =\tau^{2m}\frac{M(Y_{m+1})}{M(Y_1)}
   \prod_{k=0}^{m-1} M(-Y_k)^{-2}
 \le B_A\tau^{2m}.
\]
All the vectors are centered, so \(M(Y_1)\ge1\) and \(M(-Y_k)\ge1\). This controls the initial segment, even when \(N\) is large. At index \(N\) the three-factor bound applies and gives the rest of the estimate. If \(E_0=0\), all squared norm ratios equal \(\tau^2\). A pair with fewer active components terminates within the two sweeps discussed after \eqref{eq:rule}, and any nonzero ratios in this prefix are covered by the same bound on individual factors.

The attaining balanced example gives the sharp lower bound. In a nonterminating trajectory each sweep has \(p\) gradient updates. At most \(p-1\) factors, each bounded by \(K\), connect an intermediate gradient to the previous sweep endpoint. The endpoints provide a subsequence for the opposite inequality between root factors. A terminating trajectory has root factor zero in both index conventions.
\end{proof}

\section{Two exact examples and scope}

For \(p=1\), the balanced weighted energies at the two eigenvalues are equal. The exact sweep factor is \((\lambda_1-\lambda_0)/(\lambda_1+\lambda_0)\). This is the endpoint BB factor on the two-dimensional active subspace, consistent with the treatment of the sharp asymptotic BB factor in \cite{BBsharp}.

For \(J=\{1,2,4\}\), the values of \(c_i\) are \(8/3,2,1/3\), the balanced proportions of \(m_i\) are \(1/3,1/2,1/6\), and \(\tau=1/5\). The attaining polynomial is
\[
 P_\star(t)=1-t+t^2/5,\qquad
 P_\star(1)=1/5,\quad P_\star(2)=-1/5,\quad P_\star(4)=1/5.
\]
The root factor per individual step is \(1/\sqrt5\). The accompanying exact arithmetic script also checks a nonbalanced trajectory by comparing energy ratios and the product formula for the norm. The proof above is analytic. The script checks finite examples and does not certify the general theorem.

The sharp factor refers to one delayed complete sweep with a fixed weight and at most \(p+1\) distinct spectral nodes. Time-varying weights or reusing a filter for several sweeps changes \eqref{eq:recurrence}. With more than \(p+1\) active nodes the Vandermonde equations no longer have a one-dimensional kernel, so the same proof gives no rate on those supports. Near balance the time to leave a neighborhood defined by \(Q<A\) grows as \(\log_2(A/E_0)\), which explains why the worst linear rate can persist during a long transient even when the eventual root factor is zero.

\begin{thebibliography}{9}
\bibitem{Fletcher} R. Fletcher, A limited memory steepest descent method,
\emph{Math. Program.} 135 (2012), 413--436.
\href{https://doi.org/10.1007/s10107-011-0479-6}{doi:10.1007/s10107-011-0479-6}.

\bibitem{CurtisGuo} F. E. Curtis and W. Guo, \(R\)-linear convergence of limited memory steepest descent,
\emph{IMA J. Numer. Anal.} 38 (2018), 720--742.
\href{https://arxiv.org/abs/1610.03831}{arXiv:1610.03831}.

\bibitem{Bos} L. Bos, N. Levenberg and J. Ortega-Cerd\`a, Optimal polynomial prediction measures and extremal polynomial growth,
\emph{Constr. Approx.} 54 (2021), 431--453.
\href{https://doi.org/10.1007/s00365-020-09522-1}{doi:10.1007/s00365-020-09522-1}.

\bibitem{YangYuan} S. Yang and Y.-X. Yuan, Global and unconditional \(R\)-linear convergence of rank-compressed weighted LMSD sweeps for strictly convex quadratics,
2026. \href{https://arxiv.org/abs/2609.12114}{arXiv:2609.12114}.

\bibitem{Gu} R. Gu, Superlinear convergence of cyclic gradient methods with historical Ritz step sizes,
2026. \href{https://arxiv.org/abs/2610.32705v2}{arXiv:2610.32705v2}.

\bibitem{BBsharp} S. Yang and Y.-X. Yuan, The sharp worst-case asymptotic rate of the Barzilai--Borwein method in \(\R^d\) and Hilbert spaces,
2026. \href{https://arxiv.org/abs/2608.07839}{arXiv:2608.07839}.
\end{thebibliography}
\end{document}
